\par\noindent\textbf{(1)} 横向开路稳态条件是 \(q(E_y+\boldsymbol v\times\boldsymbol B|_y)=0\)。若电流沿 \(x\)、载流子速度为 \(v_x\)，有 \(E_y=v_xB\)（方向依坐标约定），且 \(j_x=nqv_x\)，故带符号Hall系数 \(R_H=E_y/(j_xB)=1/(nq)\) 在指定电压极性下成立。
\par\noindent\textbf{(2)} 设条宽 \(w\)，有 \(j_x=I_x/(wt)\)，电压幅值 \(V_H=|E_y|w=|R_H|I_xB/t=I_xB/(n|q|t)\)。宽度抵消，厚度不能遗漏；若使用二维面密度 \(n_s=nt\)，则 \(V_H=I_xB/(n_s|q|)\)。
\par\noindent\textbf{(3)} 横向Hall电场产生的功率密度 \(j_yE_y=0\)，因为开路 \(j_y=0\)。电流仍通过纵向电场耗散，功率 \(P=I_xV_x\ge0\)；单载流子Drude模型中纯Hall项为反对称无耗散输运，不是额外能量来源。
\par\noindent\textbf{(4)} 若电子、空穴密度为 \(n,p\)，迁移率分别 \(\mu_e,\mu_h\)，弱场两载流子Hall系数为 \(R_H=(p\mu_h^2-n\mu_e^2)/[e(p\mu_h+n\mu_e)^2]\)。因此Hall符号由迁移率平方加权决定，不能只按多数载流子的粒子数判断。
